\documentclass[12pt]{article}
\usepackage{latexsym, amscd, %graphicx, 
amssymb}
\newcommand{\SECT}[1]
		{\vspace{.3cm}
		$\!\!\!\!\!\!\!\!\!\!\!${\bf {#1}}
		\vspace{.2cm}}

\begin{document}


\begin{center}

{\large \bf Math 152, Fall, 2004, Workshop 8\par}
\vspace{.5cm}
{{\bf Honors Section} 
\par}
\end{center}

\begin{enumerate}

\item This problem is based on the fact that bounded monotonic 
sequences converge. 

\begin{enumerate}

\item Given that $1 < x < y$, find a similar relation between
${\displaystyle \frac{1}{x}}$ and ${\displaystyle \frac{1}{y}}$. 
Then show that if $1 < x < y < 4$
then ${\displaystyle 1 < 4-\frac{1}{x} 
< 4-\frac{1}{y}
< 4}$. (Note: you must show that each of these three inequalities is true). 

\item Compute five terms $(a_{1}, \dots, a_{5})$ of the recursively 
defined sequence 
$$a 1 = 1, \;\;\; a_{n+1}=4-\frac{1}{a_{n}}.$$

\item Show that $1 < a_{1} < a_{2} < 4$; then, using (a), show that the 
sequence $\{a_{n}\}$ is increasing and 
bounded. Does it converge? 

\item Find the limit of the sequence in (b) (call the limit 
$Q$, and take the $n\to \infty$ limit on both sides 
of the recursion to obtain an equation for $Q$.) 

\end{enumerate}

\item \begin{enumerate}
\item Using the integral test, show that the series 
${\displaystyle \sum_{n=2}^{\infty}\frac{1}{n(\ln n)^{2}}}$ converges

\item Let $s$ denote the sum of this series: $s = 
{\displaystyle \sum_{n=2}^{\infty}\frac{1}{n(\ln n)^{2}}}$.
If you attempted to estimate $s$ 
by computing a partial sum of the series---that is, 
by adding up the first few terms---how 
many terms would you have to include to guarantee that the error 
(the difference between 
$s$ and the sum you calculate) 
would be less than $0.05$? Explain why you are certain of this 
accuracy. (Use the inequality 
$$\sum_{n=N+1}^{\infty}f(n)\le \int_{N}^{\infty}f(x)dx.$$
You may ignore any calculator (round-off) error arising in your addition.)

\end{enumerate}

\item Consider the following three series: 
$$(\mbox{\rm i})\sum_{n=1}^{\infty}\frac{4-\sin n}{n^{2}+1},\;
(\mbox{\rm ii}) 
\sum_{n=1}^{\infty}\frac{4-\sin n}{n+1},\;
(\mbox{\rm iii}) 
\sum_{n=1}^{\infty}\frac{4-\sin n}{2^{n}+1}.$$

\begin{enumerate}

\item Use the comparison test to determine whether or not each of 
these series converges. 

\item For each series which converges, give an approximation of 
its sum, together with an error 
estimate, as follows. First calculate the sum $s_{5}$
of the first $5$ terms. Then find a number $r$ such that 
you know that the error $R_{5}=s-s_{5}$ satisfies 
$|R_{5}|\le r$. (Hint: If $0\le a_{n}\le b_{n}$ and $\sum_{n=1}^{\infty}b_{n}$
converges, then $|R_{N}|=|s-s_{N}|=|\sum_{n=N+1}^{\infty}a_{n}|\le 
\sum_{n=N=1}^{\infty}b_{n}$. If $a_{n}=f(n)$ and $f(x)$ satisfies the 
hypothesis in the integral test, then $|R_{N}|\le 
\int_{N}^{\infty}f(x)dx$. 
This is called estimating the error. In one case you will have to 
first estimate using another series, then 
estimate that error using an integral.) 

\end{enumerate}


\end{enumerate}

\end{document}

